% ---------------------------------------------------------------------------
% Appendix: Baseline Settings
% Requires: booktabs, tabularx, listings
% Suggested preamble:
%   \usepackage{booktabs,tabularx,listings}
%   \lstdefinestyle{promptstyle}{basicstyle=\ttfamily\small,
%     breaklines=true, breakautoindent=true, breakindent=20pt,
%     columns=fixed, keepspaces=true, frame=none, numbers=none,
%     aboveskip=1em, belowskip=1em}
% Matches the plain unframed listing style used by Repo2Run, App. I.2.
% ---------------------------------------------------------------------------

% Weighted wrapping column for tabularx: widths sum to the number of Y columns.
% Rename if `Y` is already taken in your preamble.
\newcolumntype{Y}[1]{>{\raggedright\arraybackslash\hsize=#1\hsize}X}

\section{Baseline Settings}
\label{app:baseline-settings}

This section records the exact configuration of each baseline agent, so that a
reader can reproduce our runs or audit where our settings depart from the
published defaults. We report three baselines: Claude Code, SWE-agent, and
SetupX.

All three run the same underlying model, \texttt{deepseek-v4-flash}, but reach
it by three different routes, because each agent speaks a different API dialect.
Claude Code requires an Anthropic-shaped endpoint and is pointed at DeepSeek's
Anthropic-compatible bridge, which maps the \texttt{claude-sonnet-*} family to
\texttt{deepseek-v4-flash}; SWE-agent resolves a LiteLLM slug and is routed
through OpenRouter with the upstream provider pinned to DeepInfra; SetupX issues
raw OpenAI-compatible HTTP and is pointed at DeepSeek's own \texttt{/v1}
endpoint. The routes are therefore not interchangeable, and we note below where
this leaves a residual difference between arms.

Each baseline is given a budget of 100 LLM calls per repository, the same
dataset and commit pin, and the same per-repository wall clock. We discuss the
limits of that budget's comparability in
Appendix~\ref{app:budget-comparability}.

% ===========================================================================
\subsection{Claude Code}
\label{app:baseline-claudecode}

Claude Code is driven as a CLI inside a per-language container, working as a
non-root user with \texttt{sudo}. The invocation:

\begin{lstlisting}[style=promptstyle]
# env: ANTHROPIC_BASE_URL=https://api.deepseek.com/anthropic
#        the CLI expands `sonnet` to claude-sonnet-5, which this bridge serves with
#        deepseek-v4-flash -- the same model the other two baselines use
#      ANTHROPIC_API_KEY=<key>     # an OAuth token is never forwarded to a third party
#      DISABLE_AUTOUPDATER=1

claude -p \
  --model sonnet \
  --max-turns 100 \
  --permission-mode bypassPermissions \
  --output-format stream-json --verbose --include-partial-messages

# The prompt is piped on stdin, not passed as an argv operand.
# No --max-budget-usd: the CLI prices every call against the Anthropic rate card
#   whatever the endpoint, so a dollar cap is not a step budget. The budget is the
#   100-turn cap together with a per-repository wall clock.
\end{lstlisting}

\paragraph{Prompt.}
Adapted from the Repo2Run instance template (\citealp{repo2run},
Appendix~I.2), changed only where our grader differs from theirs: the repository
and deliverable paths, a direct clone in place of the clone-and-copy preamble,
and the verification command. Notes 3 and 4 are imposed by the environment
rather than by the task --- the agent works as a non-root user while the image
build runs as root, and the gate invokes the system \texttt{python3}. 

\begin{lstlisting}[style=promptstyle]
We're currently setting up the environment for the following task. Here are the details:

TASK:
Set up the environment for the Python repository {full_name} (https://github.com/{full_name})
so that its test suite can be collected, and record the working setup as a Dockerfile at
/testbed/Dockerfile.gen.

INSTRUCTIONS:
Now, you'll carry out this task on your own. Your session has begun in the repository's root
directory. Use any bash commands you need. Edit and check files as needed.

The goal is to generate a Dockerfile that can successfully build and run the tests in the
repository using the command "pytest --collect-only -q".

NOTE:
1. The repository is cloned into /testbed.
2. The Dockerfile should start with the following lines (if the base image is python:3.11 and
   the repository is adamobeng/wddbfs):
```
FROM python:3.11
RUN pip install pytest
RUN git clone https://github.com/adamobeng/wddbfs.git /testbed
WORKDIR /testbed
```
3. The build runs as ROOT from an EMPTY image: drop every `sudo` prefix from the commands you
   write into the Dockerfile, and do not rely on any file from this container.
4. Install into the system Python -- do NOT create a virtualenv, since the grader runs the
   system python3.
Your task includes:
0. **Generate Dockerfile**: Create the Dockerfile at /testbed/Dockerfile.gen.
1. **Read Directory Structure**: Check the folder structure in the root directory.
2. **Check the Configuration Files**: Inspect files like "requirements.txt", "setup.py",
   "setup.cfg", "Pipfile*", etc.
3. **Determine Package Dependencies**: Handle dependencies and manage conflicting dependency
   versions.
4. **Testing**: Ensure "python -m pytest /testbed --co -q" runs without errors.
5. **Generate Dockerfile**: Write necessary installation or setup steps determined from the
   inspection. If you finish run testing successfully, you should modify the Dockerfile in
   /testbed/Dockerfile.gen.

IMPORTANT TIPS:
* Check the directory and files carefully.
* Make sure Dockerfile commands are correct.
* Use proper Dockerfile syntax and indentation.
* Test the final Dockerfile by running "python -m pytest /testbed --co -q".
\end{lstlisting}

% ===========================================================================
\subsection{SWE-agent}
\label{app:baseline-sweagent}

SWE-agent 1.1.0, run under its own \texttt{config/default.yaml} harness. The
effective configuration is reproduced below, with the runtime overrides folded
in; comments mark the values that depart from the upstream default. The two
task-carrying templates are ported from Repo2Run (\citealp{repo2run},
Appendix~I.2), with \texttt{\{\{task\_details\}\}} rebound to the variable
SWE-agent provides and the trailing shell-prompt markers dropped, since those
describe a REPL format defined by the paper's own system template, which this
arm does not use.

\begin{lstlisting}[style=promptstyle]
# env: OPENROUTER_PROVIDER=DeepInfra    # else OpenRouter may reroute between providers
agent:
  model:
    name: openrouter/deepseek/deepseek-v4-flash-0731
    per_instance_cost_limit: 2.0
    per_instance_call_limit: 100
    temperature: 0.2                    # Repo2Run App. I.2; SWE-agent default is 0.0
    completion_kwargs:
      extra_body:
        thinking:
          type: disabled                # DeepSeek v4 enables reasoning by default
  templates:
    system_template: |-
      You are a helpful assistant that can interact with a computer to solve tasks.
    instance_template: |-
      We're currently setting up the environment for the following task. Here are the
      details:

      TASK:
      {{problem_statement}}

      INSTRUCTIONS:
      Now, you'll carry out this task on your own. Your terminal session has begun in
      the repository's root directory. Use the provided commands and any bash commands
      you need. Edit and check files as needed.

      The goal is to generate a Dockerfile that can successfully build and run the
      tests in the repository using the command "pytest --collect-only -q".

      NOTE:
      1. The repository is cloned into /repo.
      2. The Dockerfile should start with the following lines (if the base image is
         python:3.10 and repository is adamobeng/wddbfs):
      ```
      FROM python:3.10
      RUN pip install pytest
      RUN git clone https://github.com/adamobeng/wddbfs.git
      RUN mkdir /repo
      RUN git config --global --add safe.directory /repo
      RUN cp -r /wddbfs/. /repo && rm -rf /wddbfs/
      RUN rm -rf /wddbfs
      ```

      Your task includes:
      0. **Generate Dockerfile**: Create a file named "Dockerfile" in the root path
         (e.g, the absolute path is /Dockerfile).
      1. **Read Directory Structure**: Check the folder structure in the root directory.
      2. **Check the Configuration Files**: Inspect files like "requirements.txt",
         "setup.py", "setup.cfg", "Pipfile*", etc.
      3. **Determine Package Dependencies**: Handle dependencies and manage conflicting
         dependency versions.
      4. **Testing**: Ensure "pytest /repo --collect-only -q" runs without errors.
      5. **Generate Dockerfile**: Write necessary installation or setup steps determined
         from the inspection. If you finish run testing successfully, you should modify
         the Dockerfile in /Dockerfile.

      IMPORTANT TIPS:
      * Check the directory and files carefully.
      * Make sure Dockerfile commands are correct.
      * Use proper Dockerfile syntax and indentation.
      * Test the final Dockerfile by running "pytest /repo --collect-only -q".
    problem_statement_template: |-
      Set up the environment for the {{language}} repository {{full_name}}
      ({{repo_url}}) so that its test suite can be collected and run, and record the
      working setup as a Dockerfile at /Dockerfile.
    next_step_template: |-
      OBSERVATION:
      {{observation}}
    next_step_no_output_template: |-
      Your command ran successfully and did not produce any output.
  tools:
    execution_timeout: 600              # upstream default: 30
    total_execution_timeout: 7200       # upstream default: 1800
    env_variables:
      PAGER: cat
      MANPAGER: cat
      LESS: -R
      PIP_PROGRESS_BAR: 'off'
      TQDM_DISABLE: '1'
      GIT_PAGER: cat
    bundles:
      - path: tools/registry
      - path: tools/edit_anthropic
      - path: tools/submit              # default.yaml also ships review_on_submit_m;
                                        # dropped, it gates submit on a git-diff review
                                        # of /repo, but the deliverable is /Dockerfile
    registry_variables:
      USE_FILEMAP: 'true'
    enable_bash_tool: true
    parse_function:
      type: function_calling
  # cache_control omitted: Anthropic-specific, and this lane is OpenAI-dialect.
  # history_processors omitted: falls through to full-history passthrough.
env:
  deployment:
    type: docker
    image: python:3.10                  # not -slim: the container needs git
    remove_container: false
    remove_images: false
  repo:
    type: local
    path: <dataset repo, checked out at the pinned base_commit>
\end{lstlisting}

% ===========================================================================
\subsection{SetupX}
\label{app:baseline-setupx}

SetupX emits no Dockerfile. It mutates a live container and records the
surviving commands to a history file, which we replay onto a pinned clone to
obtain a comparable build artifact. SetupX is configured entirely through
environment variables; it has no prompt surface we author or modify, so none is
reported here.

\begin{table}[h]
\centering\footnotesize
\caption{SetupX configuration. All settings are environment variables.}
\label{tab:setupx-config}
\begin{tabularx}{\linewidth}{@{}Y{0.95}Y{0.70}Y{1.35}@{}}
\toprule
\textbf{Variable} & \textbf{Value} & \textbf{Note} \\
\midrule
\texttt{LLM\_PROVIDER} & \texttt{openai} &
  SetupX issues raw OpenAI-compatible HTTP rather than going through
  LiteLLM. \\
\texttt{OPENAI\_BASE\_URL} & \texttt{api.deepseek.\allowbreak com/v1} &
  DeepSeek's own endpoint. \\
\texttt{OPENAI\_MODEL} & \texttt{deepseek-v4-flash} &
  The LiteLLM-style prefix is stripped, since the endpoint expects a bare model
  name. Note this is the undated slug, whereas the SWE-agent arm pins the dated
  \texttt{-0731} route; see below. \\
\texttt{DEEPSEEK\_API\_KEY} & required &
  Validated at the boundary; the run aborts rather than falling back. \\
\texttt{SETUPX\_MAX\_LLM\_\allowbreak CALLS} & 100 &
  Counted in \emph{completions}, not agent steps. See below. \\
\texttt{DOCKER\_BASE\_IMAGE} & \texttt{python:3.10} &
  SetupX's own \texttt{.env.example} default. \\
\texttt{DOCKER\_WORK\_DIR} & \texttt{/workspace} &
  The agent works in \texttt{/workspace/repo}. \\
\texttt{SETUPX\_NETWORK\_MODE} & \texttt{bridge} &
  Host networking shares the host port space, so two concurrent repositories
  collide as soon as either binds a port. \\
\texttt{SETUPX\_PHASE1\_\allowbreak TIMEOUT} & 3600\,s &
  Wall clock for the setup phase. \\
\texttt{SETUPX\_CKPT\_NS} & per-run token &
  Namespaces checkpoint images; see below. \\
\texttt{LOG\_DIR} & per-run &
  SetupX logs ${\sim}10^4$ lines per run to a file rather than stdout. Unset, it
  defaults to a shared directory where runs cannot be matched back to rows. \\
\texttt{SETUPX\_DB\_DSN} & set / unset &
  Presence enables the XPU experience store; absence pins the arm off. This is
  the single switch separating our two SetupX conditions. \\
\texttt{XPU\_ENABLED}, \texttt{XPU\_VECTOR\_ENABLED} & \texttt{1} with DSN &
  Set explicitly in both directions, so an ambient value cannot leak from one
  arm into the other. \\
\texttt{XPU\_READONLY} & \texttt{1} &
  Skips experience writeback, so concurrent repositories cannot race on
  deduplication and results do not depend on repository order. \\
\texttt{FREEZE\_TELEMETRY} & \texttt{1} & Skips telemetry writes. \\
\texttt{EMBEDDING\_API\_KEY}, \texttt{\_BASE\_URL}, \texttt{\_MODEL} &
  required with DSN &
  All three, not just the key: with only the key set, retrieval embeds against
  the wrong host, the resulting authentication failure is swallowed, and the arm
  silently degrades to no-XPU while still reporting as XPU-on. \\
\bottomrule
\end{tabularx}
\end{table}

\paragraph{Budget.}
SetupX's own \texttt{-{}-max-steps} counts agent \emph{steps}, and a single step
can cost many completions: a verification step runs the verifier's entire
sub-loop, and the second phase adds a prosecutor and a judge. Steps are
therefore not comparable with the other arms' budgets. We leave
\texttt{-{}-max-steps} unbounded and cap completions instead, at the same 100
the other two arms receive.

\paragraph{Deviations from the published implementation.}
We apply a small patch to the SetupX source, in three places.

The completion cap above is enforced at the LLM client's entry point; stock
SetupX exposes no such limit.

Checkpoint images are namespaced per run. SetupX names its snapshot images with
a single global tag, so two concurrent runs overwrite each other's snapshots and
the stale-checkpoint sweep deletes the other run's images outright --- after
which a rollback restores the wrong repository's container, or fails outright.

The read-only mode noted in the table is likewise patched in, since stock SetupX
always writes back retrieved experience.

\paragraph{Residual difference between arms.}
The SetupX arm reaches \texttt{deepseek-v4-flash} through the undated model
slug on DeepSeek's own API, while the SWE-agent arm pins the dated
\texttt{-0731} route through OpenRouter. \emph{[TODO: confirm whether DeepSeek's
direct API exposes the dated slug; if it does, pin it and delete this
paragraph.]}
